\Article{28}
Aus der Verbindung der Formeln I....IV ergibt sich die Variation des Ausdrucks $W$
\[
\begin{aligned}
\delta W={}&\int\mathrm{d}U.\delta e.\cos(4,5).[z+\alpha\alpha(\frac{1}{R}+\frac{1}{R'})]\\
&-\int\mathrm{d}P.\delta e.\cos(5,8).(\alpha\alpha\cos i-\alpha\alpha+2\beta\beta)
\end{aligned}
\]
wobei das erste Integral über alle Elemente $\mathrm{d}U$ des freien Teils der Oberfläche des Raumbereichs $s$ zu erstrecken ist, oder über die der freien Teile, falls mehrere getrennte Teile vorhanden sind. Das zweite Integral ist über alle Elemente $\mathrm{d}P$ der Linie oder der Linien zu erstrecken, die den freien Teil oder die freien Teile von den übrigen, an den Raumbereich $S$ grenzenden Teilen trennen.

Nun muss der Wert von $W$ im Gleichgewichtszustand ein Minimum sein und darf daher bei keiner unend%
\SourceLeaf{72}lich kleinen Änderung der Gestalt der Flüssigkeit, bei der das Volumen $s$ unverändert bleibt, eine negative Änderung erfahren. Bei einer solchen Gestaltänderung verschwindet also
\[
\delta s=\int\mathrm{d}U.\delta e.\cos(4,5)
\]
und man sieht leicht, dass die Oberfläche $U$ im Gleichgewichtszustand so geformt sein muss, dass in allen ihren Punkten dieses Element der Variation $\delta W$
\[
\mathrm{d}U.\delta e.\cos(4,5).\bigl[z+\alpha\alpha(\tfrac1R+\tfrac1{R'})\bigr]
\]
dem Element der Variation $\delta s$, also der Größe $\mathrm{d}U.\delta e.\cos(4,5)$, proportional ist. Gleichbedeutend damit ist
\[
z+\alpha\alpha(\tfrac1R+\tfrac1{R'})=\mathrm{Const.}
\]
Wäre diese Proportionalität nicht gegeben, so ließe sich der Wert von $W$ durch eine geeignete Änderung der Gestalt der Oberfläche $U$ offenbar vermindern, selbst wenn die Grenze $P$ unverändert bliebe. Die Gleichung gilt übrigens für die \textit{ganze} Oberfläche $U$, auch wenn sie aus mehreren getrennten Teilen besteht, sofern die Flüssigkeit selbst zusammenhängt.

Diese Gleichung bildet den ersten grundlegenden Satz in der Theorie des Gleichgewichts von Flüssigkeiten. Der berühmte \textsc{Laplace} hat ihn bereits auf einem ganz anderen Weg als dem unseren gewonnen.

Wählen wir die Ebene, für die $z$ gleich der Konstanten in der Gleichung ist und die wir normale Horizontalebene (plan de niveau) nennen können, anstelle der Ebene, von der aus die Koordinaten $z$ gemessen wurden, so gilt
\[
z=-\alpha\alpha(\tfrac1R+\tfrac1{R'})
\]
und daraus ergeben sich unmittelbar die folgenden Folgerungen.

I. Schneidet die Normalebene irgendwo die freie Oberfläche $U$, so ist die Oberfläche in jedem Punkt der Schnittlinie notwendig konkav-konvex, und der größte Radius der Konvexität ist dem größten Radius der Konkavität gleich.

II. Oberhalb der Normalebene ist die Oberfläche entweder konkav-konkav oder, wo sie konkav-konvex ist, überwiegt die konkave Krümmung die konvexe.

III. Unterhalb der Normalebene ist die Oberfläche entweder konvex-konvex oder, wo sie konkav-konvex ist, überwiegt die konvexe Krümmung die konkave.

IV. Die freie Oberfläche $U$ kann nur dann einen ebenen Teil von endlicher Ausdehnung haben, wenn er horizontal ist und mit der Normalebene zusammenfällt.

\SourceLeaf{73}\Article{29}
Gilt die soeben aufgestellte Gleichung, so reduziert sich die Variation des Werts von $W$ auf
\[
\delta W=-\int\mathrm{d}P.\delta e.\cos(5,8)(\alpha\alpha\cos i-\alpha\alpha+2\beta\beta)
\]
Führen wir also einen Winkel $A$ so ein, dass
\[
\cos A=\frac{\alpha\alpha-2\beta\beta}{\alpha\alpha}\qquad\text{oder}\qquad\sin\tfrac12 A=\frac\beta\alpha
\]
so erhalten wir
\[
\delta W=\alpha\alpha\int\mathrm{d}P.\delta e.\cos(5,8).(\cos A-\cos i)
\]
wobei die Integration über die ganze Linie $P$ erstreckt wird. Wir müssen beachten, dass der Faktor $\cos(5,8)$ gleich $\sin(5,6)$ mit positivem oder negativem Vorzeichen ist, je nachdem, ob wir uns vorstellen, dass die Flüssigkeit bei ihrer virtuellen Bewegung am Element $\mathrm{d}P$ über die Grenze $P$ hinausfließt oder sich diesseits der Grenze zurückzieht. Daraus schließen wir leicht, dass im Gleichgewichtszustand im Allgemeinen überall $i=A$ gelten muss. Wäre nämlich auf einem Teil der Linie $P$ die Beziehung $i<A$ erfüllt, so würde eine virtuelle Bewegung der ersten Art auf diesem Teil bei unverändertem restlichen Teil der Grenze $P$ offenbar eine negative Variation von $W$ hervorrufen. Ebenso ergäbe sich eine negative Variation von $W$ durch eine virtuelle Bewegung der Flüssigkeit der zweiten Art, wenn auf irgendeinem Teil der Linie $P$ die Beziehung $i>A$ bestünde. Beide Annahmen widersprechen also der Bedingung des Minimums im Gleichgewicht.

Dies ist der zweite grundlegende Satz. Er geht auch in die Untersuchungen des berühmten \textsc{Laplace} ein, doch finden wir dort keine Herleitung aus dem Prinzip der Molekularkräfte.

